% 本文件由 LaTeX 排版 Agent 根据有效 translation.json 自动生成。
% author=Athanasius; work_id=2813; title=为君士坦提乌斯辩护

\chapter{为君士坦提乌斯辩护}

% structure: p0001 source=h1 target=omitted reason=page-title-already-rendered-by-parent

1. 得知您作为最热诚的\OriginalTerm{奥古斯都}{Augustus}，多年来一直是基督徒，而且生来就敬虔，我欣然决定此刻为自己作答；——因为我将借用宗徒\Person{保禄}的话，以他作为我在您面前的辩护者，想到他既是真理的宣讲者，而您又是他话语的细心聆听者。\par

关于那些教会事务，它们已成为反对我的阴谋的借口，只需请尊驾参考许多曾为我写信的主教们的证词便已足够；\Person{乌尔撒西乌斯}和\Person{瓦伦斯}的撤销指控也足以向所有人证明，他们提出的所有指控毫无实据。其他人还能提出何等强有力的证据，能比他们书面声明的内容更强大呢？\Quote{我们说了谎，我们捏造了这些事；所有对\Person{亚大纳削}的指控全都充满虚伪。}除此之外，若您恩准垂听，还可补充一事：那些指控者在我们当面时，对\Person{马卡里乌斯}司铎提不出任何证据；但当我们不在场时，他们却随己意处理了此事。如今，首先有神圣法律，其次有我们的法律，都明确宣告这类诉讼全无效力。从这些事上，我确信尊驾作为天主和真理的热爱者，必能看出我们毫无可疑，并判定我们的对手为诬告者。\par

2. \Emph{第一项指控：挑唆\Person{君士坦斯}敌视\Person{君士坦提乌斯}。}\par

至于在尊驾面前控告我的那项诽谤指控，即关于我与最虔敬的\OriginalTerm{奥古斯都}{Augustus}——您那享有真福、永怀不忘的兄弟\Person{君士坦斯}——的通信（因为我的敌人如此报告我，并竟敢以书面断言），先前的事件已足以证明此事亦属不实。倘若指控来自另一批人，此事诚然可作详细调查，但需要有力证据，并在双方在场时公开证明：然而发明前一项指控的同一批人，又炮制了此项指控，难道不应从其中一桩的结果推断出另一桩的虚假吗？为此他们再次私下串通，以为能在我不觉时欺瞒尊驾。但他们失算了：您并未如其所愿地听信他们，反而耐心地给我机会为己辩护。您未立即采取报复，此举正合乎君王之公义，因君主的职分便是等待受害方为自己申辩。若您恩准垂听，我深信在此事上您同样会谴责那些肆无忌惮的人，他们毫不敬畏那位曾命令我们不可在君王面前妄言的天主。\par

3. \Emph{他从未单独面见\Person{君士坦斯}。}\par

但说实话，我甚至羞于为这类指控为自己辩护，我甚至以为连控告者本人都不敢在我面前提起这些。因为他完全清楚自己在说谎，而我从未如此癫狂、如此丧失理智，以致于会让人怀疑我会起这种念头。因此，若换了别人以此事质问我，我甚至不会作答，免得在辩护时，听者在那一刻暂时搁置对我的评判。但对于尊驾，我放声明确地回答，并伸出手——正如我从宗徒那里所学：\BibleQuote{我呼求天主为我的灵魂作证}\BibleRef{格后 1:23}，又如《\WorkTitle{列王纪事}》中所载（请容我同样说）：\BibleQuote{上主作证，他的受傅者作证}\BibleRef{撒上 12:5}——我从未在您兄弟、享有真福怀念的最热诚\OriginalTerm{奥古斯都}{Augustus}\Person{君士坦斯}面前，说过尊驾的坏话。我并未如他们诬告我的那样，唆使他反对您。每当我与他会面，只要他提及尊驾（他确实在\Person{塔拉苏斯}来到\Place{皮提比翁}而我住在\Place{阿奎莱亚}期间提过您），上主可以作证，我是以何等言辞论及尊驾，我巴不得天主将这话显露给您的灵魂，好让您谴责这些诽谤者的谎言。最仁慈的奥古斯都，请容忍我，并因我讲述此事而施恩准我自由陈言。您那最具基督徒精神的兄弟，绝非轻浮之人，而我也非这等品格，以至于我们竟会彼此交流这类话题，或让我在一位兄弟面前诽谤他的兄弟，或在一位皇帝面前说另一位皇帝的坏话。陛下，我没有这般癫狂，也没有忘记那神圣的话语：\BibleQuote{连心里也不可咒骂君王，在卧房里也不可咒骂富贵人：因为空中的飞鸟能传音，凡有翅膀的也会传话}\BibleRef{训 10:20}。若那些暗中对你们君王的议论尚且隐藏不住，那我怎么可能当着一位君王和诸多在场者的面诋毁您呢？因为我从未单独面见您的兄弟，他也从未私下与我交谈；我总是与当地该城的主教，以及偶然在场的其他人一起被引见。我们一同晋见，一同退出。\Place{阿奎莱亚}的主教\Person{福尔图纳提安}可以为此作证，\Person{霍西乌斯}老父也能同样证实，还有\Place{帕多瓦}的主教\Person{克里斯皮努斯}、\Place{维罗纳}的\Person{卢西鲁斯}、\Place{莱伊斯}的\Person{狄奥尼修斯}，以及\Place{坎帕尼亚}的\Person{文森提乌斯}。尽管\Place{特里尔}的\Person{马克西米努斯}和\Place{米兰}的\Person{普罗塔西乌斯}均已去世，但曾任宫廷总管的\Person{欧根尼乌斯}仍能为我作证；因为他曾立于帷幕之前，听见了我们向皇帝所求何事，及皇帝慨允如何答复我们。\par

4. \Emph{\Person{亚大纳削}的往来行踪反驳了这项指控。}\par

这确实足以证明，然而仍请容许我向您详述我的行程，这将进一步使您谴责我对手那些毫无根据的诽谤。当我离开\Place{亚历山大里亚}时，我并未前往您兄弟的营地，也未投靠其他人，而只去了\Place{罗马}；将我的案件陈明于教会（因这是我唯一关心的事），我专心于公共崇拜。我未曾给您兄弟写过信，除非是在\Person{欧瑟比}及其同党写信向他控告我，而我尚在\Place{亚历山大里亚}时被迫为自己辩护；以及后来我寄给他他命我准备的、包含\WorkTitle{圣经}的卷册时。我既为自己的行为辩护，理当向陛下陈明实情。然而过了三年，在第四年他写信给我，命我去见他（当时他在\Place{米兰}）；当我询问原因（因我对此一无所知，主为我作证），我得知有几位主教前去请求他写信给陛下，希望召开一次\OriginalTerm{公会议}{Council}。主上，请相信我，这是实情；我绝不说谎。于是我下去\Place{米兰}，并受到他极友善的接待；因为他降尊接见我，并说已致信于您，请求召开公会议。我在此城逗留时，他又召我前往\Place{高卢}（因为\Person{何西}教父正要去那里），我们好从那里前往\Place{撒狄卡}。公会议之后，他写信给我，而我仍在\Place{纳伊苏斯}，我便起身，之后住在\Place{阿奎莱亚}；在那里，陛下的书信找到了我。后来，受您已故兄弟的邀请，我再次返回\Place{高卢}，并最终来到陛下面前。\par

5. \Emph{所谓罪行，既无可能的时间，亦无可能的地点。}\par

如今，我的控告者依据他诽谤性的诬指，具体指出我在何时何地说过这些话？我当着谁的面如此疯狂，竟说出他诬陷我说的那些话？有谁愿支持这指控并为事实作证？按照\WorkTitle{圣经}的教导，他应该述说自己亲眼所见的事。但是不；他找不到这从未发生之事的证人。但我请陛下与真理一同作证，我没有说谎。我请求您——因我知道您记忆力极好——回忆您降尊接见我时，我与您的谈话：第一次在\Place{维米纳丘姆}，第二次在\Place{卡帕多细亚的凯撒勒亚}，第三次在\Place{安提约基}。我曾在您面前说过迫害我的\Person{欧瑟比}及其同党的坏话吗？我诽谤过任何伤害我的人吗？如果我在有权指责的人身上都没有指摘，我怎会如此疯狂，竟在一位皇帝面前诽谤另一位皇帝，挑拨兄弟失和？我恳求您，要么使我来到您面前，让事情得到证实，要么谴责这些诽谤，效法\Person{达味}的榜样，他说：\Quote{暗中毁谤近人的，我必予以消灭。}他们已尽其所能杀了我；因为\BibleQuote{说谎的口，杀害人的灵魂}\BibleRef{智 1:11}。但您的忍耐胜过了他们，并赐我勇气为自己辩护，好使他们这些好争闹和诽谤的人受到定罪。关于您那位至虔信、具有真福记忆的兄弟，说到这里或许就够了；因为您能凭天主所赐的智慧，从我所述不多的言语中推出许多，并认出这虚构的指控。\par

6. \Emph{第二项指控：与\Person{马格嫩提}通信。}\par

关于第二项诽谤，即我曾致信暴君（我不愿提他的名字），我恳求您以您喜欢的任何方式，并通过您所赞许的任何人，审查并核实此事。这指控的荒谬让我如此困惑，以至于我全然不知如何应对。请相信我，至虔信的君王，我多次在心里掂量这事，却不能相信有人竟如此疯狂，捏造这样的谎言。可是当\OriginalTerm{亚略派}{Arians}将这指控连同前一个指控到处散布，并吹嘘已送交您一份信件的副本时，我就更加惊讶了，我过去常常彻夜不眠，与这指控抗争，仿佛在控告者面前；我会突然放声大哭，随即跪下祈祷，悲叹流泪，切望能从您那里获得有利的听讯。现在，藉着主的恩宠，我获得了这样的听讯，却又不知该从何开始辩护；因为我每次试图开口，都因对那行为的恐惧而受阻。在您已故兄弟的事上，诽谤者确实有一个貌似可信的借口来指控；因为我曾蒙召见他，他也降尊给您这位兄弟般的仁爱写信提及我；并且他多次派人请我到他那里去，且在我到达时厚待我。但对于叛徒\Person{马格嫩提}，\BibleQuote{上主为证，他的受傅者亦为证}\BibleRef{撒上 12:5}，我不认识他，也从未与他相识。如此全然不相识的人之间，怎会有通信呢？我有什么理由要写信给这样一个人呢？如果我写信给他，我该如何起笔呢？我能说：\Quote{你杀害了那厚待我、我永难忘其恩情的人，做得好吗？} 或是：\Quote{我赞成你杀死我们的基督徒朋友和最忠信的弟兄的行为吗？} 或是：\Quote{我赞成你在\Place{罗马}屠杀那些如此友好款待我的人；比如你已故的姑母\Person{欧特罗庇亚}，她的性情正合其名，那位可敬的\Person{阿布特里}，最忠信的\Person{斯皮兰提}，以及其他许多优秀人士吗？}\par

7. \Emph{此项指控，全不可信且荒谬。}\par

我的控告者竟怀疑我有这种事，岂非纯粹疯狂？我再次请问，是什么能诱使我去信任这个人？我在他的品格中看到了哪一点可以让我倚靠？他杀害了自己的主人；他对朋友们背信弃义；他违背了自己的誓言；他违背天主的法律去请教下毒者和术士，从而亵渎了\Person{天主}。我又凭着什么良心，竟向这样一个人致意——他的疯狂与残暴不仅折磨了我，也折磨了我周围的世界？确实，他对我所为，令我对他\Quote{感激涕零}：当你已故的兄弟以圣洁的礼物装满我们的\OriginalTerm{教会}{Church}时，他却将他杀害。因为这个恶徒，看到这些礼物仍无动于衷，也不敬畏那在\OriginalTerm{圣洗}{baptism}中赐给他的\OriginalTerm{恩宠}{divine grace}，反而像一个受诅咒的魔鬼附身者，对他暴怒，直至你有福的兄弟在他手下惨遭\OriginalTerm{殉道}{martyrdom}；而他从此成了一个像\Person{加音}一样的罪犯，被赶逐流离，\BibleQuote{呻吟颤抖}，最终效法\Person{犹达斯}之死，自戕而亡，从而在将来的审判中，给自己带来双倍的惩罚。\par

8. \Emph{反驳}\par

与这样一个人，这个诬告者竟以为我曾有友谊之交，或者更确切地说，他并非果真如此认为，而是像一个仇敌般编造了一个难以置信的谎言：因为他完全清楚自己在撒谎。我但愿，不管他是谁，他能在此，我便可以凭着\OriginalTerm{真理}{Truth}本身（因为我们\OriginalTerm{基督徒}{Christian}无论说什么，如同在\Person{天主}面前，都视为起誓）去质问他；我说，我会问他这个问题：在已故的\Person{君士坦斯}的福祉中，我们谁更为喜乐？谁为他祈祷得更恳切？前述指控的事实证明了这一点；其实，情况如何，人人了然。然而，尽管他自己完全清楚，凡对已故的\Person{君士坦斯}抱有如此心意且真正爱他的人，不可能与他的仇敌为友，但我担心，他因对他怀有不同于我的情感，便将他心怀的仇恨之情错归于我。\par

9. \Emph{\Person{亚大纳削}不可能写信给一个甚至不认识他的人}\par

就我自己而言，我对这事的滔天罪行如此震惊，以致全然不知该说些什么来为自己辩护。我只能宣明：若此事于我稍有嫌疑，我愿判自己万死。陛下，您作为真理的爱好者，我向您坦荡申诉。我恳求您，如我之前所说，调查此事，特别是用那些曾由他派来作为使者觐见您的人作证。这些人是\Person{萨尔瓦提乌斯}主教、\Person{马克西穆斯}主教以及其他人，还有\Person{克莱孟提乌斯}和\Person{瓦伦斯}。我恳求您，问问他们，是否曾带信给我。若他们带过，这便给了我写信给他的缘由。若他未曾写信给我，甚至根本不认识我，我怎能写信给一个不认识的人呢？问问他们，当我见到\Person{克莱孟提乌斯}及其同伴，并提起您那满怀福忆的兄弟时，我难道没有，按照\OriginalTerm{圣经}{Scripture}的言辞，用眼泪沾湿我的衣裳吗？因我想起他的亲切性情和\OriginalTerm{基督徒}{Christian}心肠。从他们那里，请听我得知那野兽的残忍，并发现\Person{瓦伦斯}及其同伙取道\Place{利比亚}时，是何等焦虑，唯恐他也会企图渡海，像强盗一样杀害那些在爱中忆念已故君王的人——我自认在这群人当中不落人后。\par

10. \Emph{他对\Person{君士坦提乌斯}及其兄弟的忠诚}\par

既然对他们这种意图有所忧虑，我岂不更应为陛下祈祷吗？难道我会对杀害他的凶手怀有感情，而对为他复仇的您——他的兄弟，怀有厌恶吗？我若记得他的罪行，岂能忘记您那份善意——承您垂允，曾以书信确保您在我心中，在您那位满怀福忆的兄弟去世后，仍将与他在世时一样？我怎忍直视那凶手呢？当我为您的安全祈祷时，难道不会想到那位蒙福的君王在注视着我吗？因为兄弟天生是彼此的明镜。因此，在他（\Person{君士坦斯}）身上看见您，我绝未在他面前毁谤您；在您身上看见他，我也绝未写信给他的仇敌，而不为您祈祷。为此，我的证人有：首先是\OriginalTerm{主}{Lord}，祂俯听了我，并将您祖先的全部王国赐予了您；其次是当时在场的人们：\Person{费利奇西穆斯}，那位\Place{埃及}公爵；\Person{鲁菲努斯}和\Person{斯德望}，前者是总收税官，后者是当地长官；\Person{阿斯忒里乌斯}伯爵，以及宫廷总管\Person{巴拉狄乌斯}；还有公务官员\Person{安提约古}和\Person{埃瓦格里乌斯}。我只需说：\Quote{让我们为最虔诚的皇帝、\Person{君士坦提乌斯}·奥古斯都的安全祈祷，}所有民众便立即同声呼喊：\Quote{哦，\Person{基督}，援助\Person{君士坦提乌斯}！}他们如此持续祈祷了一段时辰。\par

11. \Emph{就所谓信件向控告者提出挑战}\par

现在我已经呼求天主，以及祂的圣言，我们的主耶稣基督，\OriginalTerm{独生子}{Only-begotten Son}，为我作证，我从未给那人写过信，也未收到过他的信。至于我的控告者，请容我就这项指控向他提几个简短问题。他是如何得知此事？他会说他获得了信件的副本吗？因为\OriginalTerm{亚略派}{Arians}曾竭力证实此事。而首先，即使他能出示与我的笔迹相似的文字，事情也并非确凿；因为存在着伪造者，他们甚至常常模仿你们这些皇帝的手迹。笔迹相似并不证明信件的真伪，除非我的常用\OriginalTerm{秘书}{amanuensis}为之作证。那么我要再问我的控告者，是谁向你们提供了这些副本？它们又是从何而来？我有我的书记员，而他也有他的仆人，这些仆人从送信人那里接收信件，并交到他手中。我的助手可以找到；敬请传唤其他人（他们极可能尚在世），并就这些信件进行调查。请彻查此事，如同真理与你的宝座相伴。她是君王的保障，尤其是基督徒君王的保障；有她同在，你必能最稳固地统治，因为圣经上说：\BibleQuote{仁爱与忠信保全君王，他的宝座因仁义而稳固}\BibleRef{箴 20:28}。智慧的\Person{则鲁巴贝耳}曾因阐明真理的力量而胜过其他人，全体民众遂高呼：\BibleQuote{真理伟大，远超一切}\BibleRef{厄三 4:41}。\par

12. \Emph{真理乃宝座的保障。}\par

倘若我是在别的任何人面前被指控，我早该向你的虔诚上诉了；就如昔日那位\OriginalTerm{宗徒}{Apostle}向\Person{凯撒}上诉，从而挫败了他的仇敌针对他的阴谋。然而，如今他们竟敢向你们提出指控，我又能离开你向谁上诉呢？唯有向那位曾说过 \BibleQuote{我就是真理}\BibleRef{若 14:6} 者的父上诉，求祂使你的心转向仁慈：—\par

全能的主，永恒的君王，我们的主耶稣基督的父啊，祢曾藉着祢的\OriginalTerm{圣言}{Word}将这王国赐给祢的仆人\Person{君士坦提乌斯}；求祢光照他的心，使他识破攻击我的虚谎，得以悦纳我这辩护；并向万民昭示，他的耳专注倾听真理，正如经上所载：\BibleQuote{只有正义的嘴唇才蒙君王悦纳}。因为祢曾藉\Person{撒罗满}的口说，这样，王国的宝座将得以坚立。\par

因此，至少请你调查此事，让那些控告者明白，你的愿望是了解真理；且看，他们的神情是否会显露他们的虚谎；因为面容是良心的试金石，正如经上记载：\BibleQuote{愉快的心使面容焕发，心灵忧伤使精神颓丧}\BibleRef{箴 15:13}。就这样，那些合谋陷害\Person{若瑟}的人，因自己的良心而被定罪；\Person{拉班}对\Person{雅各伯}的诡计，也在面容上显露出来。同样，你也可看到这些人疑虑惊惶，因为他们逃遁躲避；而我们则坦然自辩。我们之间的争讼并非关乎尘世财富，而是关乎教会的荣耀。被石头击伤的人，会求医治疗；然而诬蔑的打击比石头更锋利；正如\Person{撒罗满}所说：\BibleQuote{作假见证的人，无异槌、刀和利箭}\BibleRef{箴 25:18}，其创伤唯有真理才能治愈；若真理被蔑视，创伤便会日益恶化。\par

13. \Emph{指控基于伪造。}\par

正是此事使得各地教会陷入如此混乱；因为种种借口被编造出来，一些德高望重、年长的\OriginalTerm{主教}{Bishops}，仅因与我保持共融而被放逐。若事情仅止于此，透过您仁慈的干预，我们的前景尚属可期。然而，为防止这邪恶蔓延，让真理在您面前得胜；不要使各教会都蒙受猜疑，仿佛基督徒、甚至主教们竟会如此密谋和通信。若您不愿彻查此事，那么理应相信我们自辩者，而非那些\OriginalTerm{诬告者}{calumniators}。他们犹如仇敌，专行邪恶；我们则竭力为己案申辩，向您呈上证词。我确感诧异，为何我们怀着敬畏与虔诚向您陈情，他们却厚颜无耻，竟敢在皇帝面前撒谎。但我恳求您，看在真理的份上，并如经上所载，\BibleQuote{细加查考}，在我面前，查明他们根据什么断言这些事，这些信件从何而来。然而，我的仆人皆无罪行可被证实，他的人也无人能说出信件的来源；因为这些信件是伪造的。也许，最好不要再追问下去了。他们不希望追问，免得那写信者必定被查出。因为除了那些诬告者，再无他人知道他是谁。\par

14. \Emph{第三项指控：使用\OriginalTerm{未祝圣的教堂}{undedicated Church}。}\par

因他们在\\Place{大教堂}的事上控告我，说那教堂尚未完工就在其中举行了\\OriginalTerm{共融礼仪}{communion}，我要就此事向陛下作答；因为那些敌视我的人迫使我这样做。我承认此事确有发生；因为正如我迄今所言，我从未说过谎言，现在我也不会否认此事。但事实远非他们所描述的那样。请容我向您禀明，至虔敬的奥古斯都，我们并没有举行任何祝圣典礼（未经陛下下令，这样做当然是不合法的），我们的行为也并非出于预谋。当时没有邀请任何主教或其他圣职人员参与我们的集会；因为建筑尚有许多未竟之处。而且这次聚会并非事先通知，好让他们有机会控告我们。每个人都知道事情是如何发生的；但请以您惯常的公正与耐心听我道来。那是复活节庆日，聚集的群众极其众多，如此盛况，正是基督信徒的君王愿意在所有城市中见到的。当时各教堂都显得太小，容纳不下众人，于是民众中起了不小的骚动，他们要求获准在\\Place{大教堂}内聚会，在那里他们可以一起为陛下的安康献上祈祷。他们也这样做了；因为尽管我劝导他们稍安毋躁，并到其他教堂去举行礼仪，无论自己多么不便，他们却不听我的劝告；他们已准备出城，到旷野露天的地方聚会，认为宁可忍受旅途劳顿，也不愿在如此拥挤不适的情况下过节。\par

15. \\Emph{因场所不足，且有成例可证。}\par

陛下，请相信我，并让真理也为我作证，我在此声明：在四旬期的集会中，由于场所狭窄和民众聚集极多，许多孩童、不少青年、极多的老年妇女，以及好几位青年男子，因人群拥挤而大为受苦，以至于不得不被人抬回家中；虽然赖\\OriginalTerm{天主眷顾}{Providence}，无人死亡。但众人都抱怨不已，要求使用\\Place{大教堂}。如果节日前几天就已如此拥挤，那么节日期间又会怎样呢？当然情况会糟糕得多。因此，将民众的喜悦化为忧伤，将他们的欢快化为悲愁，使节庆变为哀悼的时节，这实非所宜。\par

我更有理由如此行，因为我有教父们的先例可循。\\Person{圣亚历山大}，当其他场所过于狭窄，而他正在建造当时被视为很大的\\Place{提奥纳教堂}时，由于信徒人数众多，他既在那里举行集会，同时又继续建造工程。我在\\Place{特里尔}和\\Place{阿奎莱亚}也见过同样的事，在这两处，工程进行期间，因人数之故，在节日时他们也在该处聚集，而从未有人以这种方式控告他们。不但如此，您那永怀不忘的兄弟也曾亲身在场，那时在\\Place{阿奎莱亚}便是在这种情形下举行了\\OriginalTerm{共融礼仪}{communion}。我也效法了此一作法。并没有祝圣典礼，仅仅是一场祈祷礼仪。至少我确信，您作为热爱天主的人，定会赞许民众的热心，并恕宥我不愿阻止如此众多信众的祈祷。\par

16. \\Emph{一起祈祷胜于分开祈祷。}\par

然而，在这里我要反问我的控告者：民众应在哪里祈祷呢？在旷野中，还是在为祈祷之目的而正在建造的地方？民众应答\\Quote{阿们}，哪里才是合宜且虔诚的呢？是在旷野，还是在已被称为上主之殿的地方？至虔敬的君王，您愿意您的子民在哪里伸开双手为您祈求呢？愿意在希腊人路过时会驻足聆听的地方，还是在以您命名的、自奠基以来便久已称为上主之殿的地方？我确信您会偏爱您自己的地方；因为您在微笑，这告诉我如此。控告者却说：\\Quote{应该在教堂里。}但正如我之前所说，所有教堂都太小、太狭窄，容纳不下如此之众。再者，哪种方式最适宜他们祈祷呢？是让他们分作数处、分成小群，冒着拥挤会众的危险聚会呢？还是既然已有一个能容纳所有人的地方，就在其中共同聚集，并如同以同一声音完全和谐地诵念呢？后者更好，因为这显示出民众的同心合意：天主也必乐于垂听这样的祈祷。因为按我们救主亲许的诺言（\\BibleRef{玛 18:19}），若两人同心合意，无论求什么，必为他们成就；那么，当如此庞大的会众同声向天主高呼\\Quote{阿们}时，又当如何呢？事实上，谁见了不感到惊奇？谁见了如此众多会众聚集一处，不对您称福呢？民众自己，因以往惯于分散各处聚会，现在彼此相见，是多么欢喜！这件事给所有人带来喜乐，惟独使那诽谤者恼怒。\par

17. \\Emph{在建筑物内祈祷胜于在旷野。}\par

现在，我还要应对控告者对我提出的其余且仅剩的一项反对。他说，建筑尚未完工，不应在那里举行祈祷。但主曾说过：\\BibleQuote{但你祈祷时，要进入你的内室，关上门}\\BibleRef{玛 6:6}。那么，控告者将如何回答？更确切地说，所有明智且真正的基督徒将说些什么？让陛下询问这类人的意见：因为经上关于前者写道：\\Quote{愚昧人将说愚昧话}；而关于后者则写道：\\BibleQuote{你当向一切有智慧的人求教}\\BibleRef{多 4:18}。当教会太过狭小，而民众如此众多，并且渴望走出去到旷野中去时，我该做什么呢？旷野没有门，凡愿意的人都可以通过，但主的殿宇却由墙壁和门围护，将虔敬者与亵渎者区分开来。那么，难道每一个明智的人，连同你虔敬的心，陛下，不会偏向后一个地方吗？因为他们知道，在那里祈祷是合法地奉献的，而在旷野却带有不合规矩的嫌疑。除非确实没有任何适宜祈祷的地方，且敬拜者只居住在旷野中，如\\Person{以色列}人的情况那样；虽然在他们建造了\\OriginalTerm{帐幕}{tabernacle}后，从此便有了一个专为祈祷而设的地方。啊！\\Person{基督}，主及万王的真王，天主的独生子，圣父的\\OriginalTerm{圣言}{Word}与\\OriginalTerm{智慧}{Wisdom}，我受指控，乃因民众祈求你的恩宠，并藉着你恳求你的父——那位在万有之上的天主，拯救你的仆人至虔敬的\\Person{君士坦提乌斯}。但感谢你的仁慈，我是因此而受责备，并且因为遵守你的法律。倘若我们绕过了皇帝正在建造的地方，而走向旷野去祈祷，那责备定会更重，指控也更为真实。那时，控告者将如何发泄他的愚妄！他会以何等貌似的理由说：\\Quote{他轻视了你正在建造的地方；他不赞同你的工程；他嘲弄地经过那地；他指出旷野来弥补空间的不足；在人民想要献上祈祷时，他加以阻止。}这正是他想说的话，也是他设法找寻机会要说的话；他没有找到机会，便恼怒不已，于是立刻捏造一项控告来攻击我。倘若他能说出这些话，他定会用羞耻来击败我；正如现在他伤害我，效法控告者的行径，并伺机攻击那些祈祷的人。他就是这样将他所知的\\BibleBook{}{达尼尔}的事迹\\BibleRef{达 6:11}扭曲为邪恶的目的。但他受骗了；因为他无知地以为巴比伦的习俗在你们中间流行，却不知道你是真福\\Person{达尼尔}的朋友，敬拜同一位天主，并不禁止，反而希望所有人都祈祷，深知所有人的祈祷都为求你继续在永久的和平与安全中掌权。\par

18. \\Emph{先有祈祷，并不妨碍后来的奉献礼}\par

这就是我须就控告者而言要抱怨的。但愿你，至虔敬的奥古斯都，在未来的许多年中长寿，并举行此教堂的奉献礼。诚然，众人为你的安全而献上的祈祷，绝不会妨碍这场庆典。愿这些无学之人停止如此歪曲，倒不如让他们从教父们的榜样中学习；并让他们研读\\WorkTitle{圣经}。或者更好，让他们从你学习——你对这类历史如此精通——那位司祭约匝达克的儿子\\Person{耶叔亚}，他的弟兄，那明智的沙耳提耳的儿子\\Person{所罗巴伯}，以及司祭兼法律经师\\Person{厄斯德拉}，当圣殿在被掳期后建造之时，帐棚节临近了（这是在\\Person{以色列}中举行盛大集会与祈祷的节期与时日），他们同心合意地聚集百姓在第一道门内的朝东大院里，为天主准备了祭台，在那里奉献了礼物，并举行了节庆。此后，他们在安息日和月朔时，还将祭品带到那里，百姓也献上祈祷。然而\\WorkTitle{圣经}明说，当这些事进行时，天主的殿宇尚未建成；反倒在他们如此祈祷时，殿宇的工程却在进展。这样，他们的祈祷并未因期待奉献礼而延迟，而奉献礼也未因祈祷的集会而受阻。百姓持续如此祈祷；当殿宇完全完工时，他们举行了奉献礼，并为此带去了礼物，众人都因工程的完成而守了庆典。真福\\Person{亚历山大}和其他教父们也这样做了。他们继续聚集他们的百姓，完工时他们感谢主，并举行了奉献礼。这也正是你当作的，啊，君王，你这最用心查考的人。这座地方已经预备好了，已通过在其中献上的祈祷而得到圣化，只待你虔敬之心亲临。它的全美仅缺此一项。那么请补足这一缺失，在那里向主献上祈祷，就是你为他建造这殿宇的那一位。你能如此行，正是所有人的祈祷。\par

19. \\Emph{第四项指控：违抗帝国命令}\par

如今，若您乐意，请让我们考虑那剩余的指控，并容我同样作出答辩。他们竟敢于指控我抗拒您的命令，并拒绝离开我的教会。我实在诧异他们竟不倦于散布诽谤之词；至于我，却尚未倦于答辩，我反倒欢喜如此；因为我的答辩越充分，他们就越必须被定罪。我没有抗拒您的虔敬的命令，天主不许！我岂是那等会抗拒本城\OriginalTerm{财务官}{Quæstor}的人，更不用说如此伟大的君王了。关于此事，我无需多言，因全城都会为我作证。然而，请容我再从头叙述当时的情况；因为当您听到后，我确信您会因我敌人的狂妄而震惊。\Person{蒙塔努斯}，宫廷的官员，前来并带给我一封信，该信自称是对我的一封信的回信，我在那封信中请求前往\Place{意大利}，为的是获得补给以弥补我认为我们众教会的状况中的不足。现在，我愿感谢您的虔敬，您屈尊应允了我的请求，假定我曾写信给您，并为我安排了旅行的供应，使我得以顺利完成，不受烦扰。但在此，我再次对那些在您耳边说谎的人感到诧异，他们竟不畏惧，因为说谎属于魔鬼，说谎者与那位说\BibleQuote{我是真理}\BibleRef{若 14:6}的毫不相干。因我从未写信给您，我的指控者也无法找出任何这样的信来；虽然我本应每日写信，若能因此得见您仁慈的容颜，然而，抛弃众教会既非虔诚之举，搅扰您的虔敬亦非正当，尤其因为您虽不待我们当面提出，却乐意应允我们为教会所呈的请求。如今，请您恩准我诵读\Person{蒙塔努斯}命令我做的事。内容如下。***\par

20. \Emph{他违命的经过。}\par

如今我再问，我的指控者们又是从何处得到这封信的呢？我愿从他们那里得知，是谁将此信交在他们手中的？请您吩咐他们回答。由此您可看出，他们伪造了这封信，正如他们散布先前那封针对我而发出、涉及恶名昭彰的\Person{马格嫩提乌斯}的信一样。他们在此事上也被证实有罪，那么下次他们又将借何口实押我来受审呢？他们唯一关切的就是将一切陷入混乱和无序；我看出他们正为此目的而大发热心。也许他们以为，通过反复指控，最终能激怒您来反对我。但您应当远离这样的人，并憎恶他们；因为他们是怎样的人，也就以为听他们的人也是怎样的人；他们以为自己的诽谤甚至能在您面前得逞。昔日\Person{多厄格}\BibleRef{撒上 22:9}的指控胜过了天主的司祭们，但那是邪恶的\Person{撒乌耳}听信了他。\Person{依则贝耳}能以她的诬告伤害至为虔诚的\Person{纳波特}\BibleRef{列上 21:10}，但那时是邪恶且背教的\Person{阿哈布}听信了她。然而至圣的\Person{达味}，您当效法的榜样（正如众人祈求您能如此），并不恩待这样的人，反惯于转面不顾，如同躲避狂犬。他说：\BibleQuote{凡暗中毁谤邻人的，我必灭绝。}因他谨守诫命，其中说：\BibleQuote{你不可接受虚假的报告}\BibleRef{出 23:1}。而这些人…在您眼前所报的正是虚妄。您如同\Person{撒罗满}，曾向主要求（您当相信已蒙应允），求祂乐意使虚谎的言语远离您\BibleRef{箴 30:8}。\par

21. 那么既然这封信源于一个虚假的故事，并且没有包含要我去见您的命令，我便断定，您的虔敬并不希望我去。因为在信中您并没有给我绝对的命令，而只像是答复我请求准许整饬那些看似不足之事的信函，因此对我而言（尽管无人告诉我），显然我所收到的信并非表达您仁慈的心意。众人都知道，我也书面声明，正如\Person{蒙塔努斯}也清楚，我并不拒绝前往，只是我认为利用我曾写信向您请求这一恩典的假设是不合宜的，也担心诬告者会以此为口实，说我搅扰您的虔敬。尽管如此，我已做好准备，\Person{蒙塔努斯}也知道，以便若您屈尊给我写信，我可以立即离家，欣然应命；因为我还没有疯狂到抗拒您如此命令的地步。那么，事实上您的虔敬并未给我写信，我如何能抗拒我从未收到的命令呢？他们怎能说我拒不从命，而实际上无人给我下达命令呢？这难道不又是敌人捏造出来的、装作发生过的事吗？我担心，即使现在我正在进行这番自我辩护时，他们仍可能指控我做了从未得到您许可的事。我的行为如此轻易地就被他们作为指控的把柄，他们如此迫不及待地宣泄其诽谤，不顾圣经所言：\BibleQuote{不要喜爱毁谤他人，免得你被剪除。}\par

22. \Emph{\Person{狄奥格内斯}与\Person{西里阿诺斯}的到来。}\par

过了二十六个月，\Person{蒙塔努斯}离去之后，书记官\Person{狄奥吉尼}来了；但他没有带给我任何书信，我们没有见面，他也没有传达任何来自你的命令。再者，当将军\Person{叙利亚努斯}进入\Place{亚历山大里亚}时，看到\OriginalTerm{亚略派}{Arians}散布某些谣言，宣称事情现在将如他们所愿，我便询问他是否带来了关于他们这些说法的任何信件。我承认我请求过含有你的命令的信件。当他说他没有带来时，我请求\Person{叙利亚努斯}本人，或埃及总督\Person{马克西穆斯}，就此事项写信给我。我提出这个请求，是因为陛下曾写信给我，希望我不让自己受任何人惊扰，也不要理会那些想恐吓我的人，而要我继续毫无惧怕地居住在教会里。带来这封信的是宫廷总管\Person{帕拉第乌斯}和前亚美尼亚公爵\Person{阿斯忒里乌斯}。请允许我读一读它的副本。内容如下：\par

23. \Emph{书信副本}如下：\par

\Person{君士坦提乌斯·维克托·奥古斯都}致\Person{亚大纳削}：明智者，你并非不知道，我如何不断地祈祷，愿我已故兄弟\Person{君士坦斯}在他一切事业中获得成功；而你的智慧也很容易判断，当我得知他被歹人的阴谋害死时，我是多么痛苦。现在因为某些人正在这时试图惊吓你，把那一可悲的惨剧摆在你眼前，我认为宜向你这可尊敬的人发出这封信，劝勉你，要如主教所当为的，教导人民遵行正统信仰，并按你的习惯，和他们一起专心祈祷。因为这是合乎我们的愿望的；我们的愿望是，你应当随时在你本处作主教。\par

另笔：——愿\OriginalTerm{天主眷顾}{divine Providence}保全你，亲爱的神父，长寿多年。\par

24. \Emph{亚大纳削为何不服从谕旨}\par

关于这封信，我的反对者们与地方官员商议。我既然收到了它，要求他们出示信件，并拒绝理会那些纯粹的借口，这难道不合理吗？他们未能向我出示陛下的谕令，这不正是直接违背了您给我的指示的要旨吗？我因此，看出他们拿不出任何您的信件，认为他们不可能仅仅获得口头传达，尤其是因为陛下的信曾嘱咐我不要听信这类人。那么，至虔诚的奥古斯都啊，我的做法是正当的：我既凭借您的信函回到本国，就应当仅凭您的命令离开；这样日后我便不致背负离弃教会的罪名，而是因接到您的命令，才有理由退隐。所有民众同司铎们一道去见\Person{叙利亚努斯}，城里至少大部分人也都一同前往，为我提出了这个请求。埃及总督\Person{马克西穆斯}也在场：他们的请求是，要么他向我书面宣告您的意愿，要么停止搅扰教会，同时民众自己正就此事项派遣代表前往您那里。当他们坚持这个请求时，\Person{叙利亚努斯}最终认识到它的合理性，并同意了，还以您的安全起誓（\Person{希拉里}在场并为此作证），说他将终止这场骚乱，并将此案件呈报陛下。公爵的卫兵以及埃及总督的卫兵都知道这是真实的；该城的\Person{普里塔尼斯}也记得这些话；因此您会看出，不论是我还是其他任何人，都没有抗拒您的谕令。\par

25. \Emph{叙利亚努斯闯进教堂}\par

所有人都要求出示陛下的信函。因为尽管君王的空口之言与其书面谕令有同等的分量和权威，尤其是传达者若敢于书面肯定那是赐给他的；然而，他们既没有公开声明已收到任何谕令，也没有依照请求向我出具书面保证，反而完全是按自己的权威行事；我承认，我冒昧地说，我对他们起了疑心。因为他们身边有许多\OriginalTerm{亚略派}{Arians}的人，是他们的座上客和谋士；他们虽未公开图谋什么，却准备用诡计和奸诈攻击我。他们的行动全不像是奉行王室谕令，反而，正如其行为显露的，是受敌人唆使。这使我更迫切地要求他们出示您的信函，因为他们的所为和计划都令人生疑；也因为在我凭借您那么多信函的权威进入教会之后，若没有这样的许可就退离，是很不体面的。可是\Person{叙利亚努斯}作出承诺之后，所有民众都满怀喜乐和安心，聚集在教会里。但二十三天之后，他带着兵丁闯进教堂，当时我们正进行惯常的事奉，正如那些当时进去的人所见证的；因为那是一个守夜祈祷，预备次日的共融。那夜所发生的事，正如\OriginalTerm{亚略派}{Arians}所渴望并早已预言要加在我们身上的。因为将军把他们带了来；他们是这次攻击的煽动者和策划者。至虔诚的奥古斯都啊，我这并非不可信的故事；因为它不是暗中进行的，而是到处哄传的。因此，当我看到攻击开始时，我先劝勉民众退去，然后自己也随之退去，天主将我隐藏并引导我，当时与我在一起的人可为此作证。从那时起，我独自隐居，不过我完全自信能为我的行为答辩，首先在天主面前，其次在陛下面前；我并没有逃跑离弃我的子民，反倒能指出将军对我们的攻击，作为迫害的证据。他的行为使所有人极其惊骇；因为他要么不应作出承诺，要么作出之后不应违背。\par

26. \Emph{亚大纳削在事发时的做法}\par

现在他们为何要设下这计谋陷害我，并背信弃义地埋伏捉拿我呢？其实他们本可以通过书面命令来强制执行。皇帝的命令通常会使执行者胆大包天；但他们鬼祟行事，反使人更怀疑他们并未领受任何命令。我要求了什么十分荒谬的事吗？让陛下明察。谁不会说，一位主教作此请求是合理的？您知道，因为您读过圣经，一位主教遗弃他的教会、疏于照管天主的羊群，是何等重大的冒犯。因为牧者不在，便给豺狼机会攻击羊群。这正是亚略异端者及所有其他异端者渴望的，乘我不在之机，诱使民众陷入不虔。假如我逃走了，在真正的主教们面前，我怎能辩护？更确切地说，在将祂的羊群托付给我的那位面前？那位审判全地的，万有的真正君王，我们的主耶稣基督，天主子。人人岂不要合理地指责我疏忽了我的子民吗？虔诚的陛下岂不会责备我，并公正地诘问：\Quote{您在接到我们的书信授权返回之后，为何又未经授权而离开，抛弃您的子民呢？}在审判的日子，子民们岂不会理所当然地把我的这疏忽归咎于我，并说：\Quote{那监管我们的人逃跑了，我们被疏忽了，没有人提醒我们尽本分。}当他们这样说时，我能回答什么呢？昔日\Person{厄则克耳}就曾如此控诉昔日的牧者；而真福宗徒\Person{保禄}深知于此，便藉他的门弟子叮嘱我们每一个人说：\BibleQuote{不要疏忽你心内的神恩，即从前藉长老团的覆手赐给你的神恩。}\BibleRef{弟前 4:14}我因畏惧这话，不愿逃走，而愿领受您的命令，如果这确是陛下您的旨意。但我从未得到我所合理请求的，如今却在您面前被人诬告；因为我没有抗拒过陛下您的任何命令；现在我也不会试图返回\Place{亚历山大里亚}，直到陛下您愿意之时。我预先说明这一点，免得那些诽谤者又以此作借口来指控我。\par

\Emph{\Person{亚大纳削}离开\Place{亚历山大里亚}前往\Person{君士坦提乌斯}处，但被主教们遭放逐的消息阻止。}\par

我注意到了这些事，没有对自己定罪，反而带着这番辩词，匆匆前来见虔诚的陛下，因为知道您的良善，记着您信实的许诺，并且深信正如天主的\WorkTitle{箴言}所载：\BibleQuote{正义的言辞，为仁君所悦纳}。但当我已经踏上旅途，穿越了沙漠之后，忽然有消息传到我耳中，起初我以为难以置信，但后来却证实是真的。四处都风传\Place{罗马}的主教\Person{利贝里乌斯}、\Place{西班牙}的伟大的\Person{贺西乌斯}、\Place{高卢}的\Person{保利诺}、\Place{意大利}的\Person{狄奥尼修斯}与\Person{欧瑟比乌斯}、\Place{撒丁岛}的\Person{路济弗尔}，以及其他若干主教、长老和执事，因拒绝签署对我的定罪而被放逐。这些人已被放逐；而\Place{卡普亚}的\Person{文森修}、\Place{阿奎莱亚}的\Person{福尔图纳提安}、\Place{得撒洛尼}的\Person{赫雷米}，以及所有\Place{西方}的主教，都受到了非比寻常的强制手段，甚至遭受极端暴力和严重伤害，直至被迫承诺不与我共融。正当我为这些消息震惊困惑之际，瞧，又有另一个报告传到，是关于\Place{埃及}和\Place{利比亚}的：说将近九十位主教受到迫害，他们的教会已交给亚略异端的信奉者；十六人被流放，其余人中，有的逃亡，有的被迫假装顺从。据说那里的迫害如此酷烈，以致在\Place{亚历山大里亚}，当弟兄们在复活期和主日，于公墓附近的一处旷野中祈祷时，将领率领超过三千名士兵，手持武器、出鞘的剑和长矛，突然袭击他们；随之而来的便是对那些仅仅在向天主祈祷的妇女和儿童所施加的暴行，其恶劣程度正如对毫无防备之人发起攻击后所能预料的一样。此时详细叙述或许不合时宜，以免单单提起这些暴行就使我们人人落泪。但他们竟如此残忍，以致贞女被剥去衣服，而那些因受击打而死的人，尸体甚至未立即交付埋葬，反被抛弃给狗，直到他们的亲属冒着极大危险，暗中前来将尸体偷走，并且需费尽心力，不使任何人知晓。\par

\Emph{关于\Person{乔治}闯入的消息。}\par

他们其余的行为，或许会令人难以置信，因其极端的残忍而令所有人震惊。然而，有必要提及这些事，好让您基督徒的热忱与虔诚能够察觉，他们对我们所构陷的毁谤与诬蔑，其目的无非是要将我们逐出教会，并以他们自己的不虔敬取而代之。因为，当那些合法的主教，即年高德劭之士，有的被放逐，有的被迫逃亡之后，\OriginalTerm{阿里乌斯派}{Arians}便立即委派\OriginalTerm{异教徒}{heathens}和\OriginalTerm{慕道者}{catechumens}，即那些在元老院中居首位的人和众所周知的富人来宣讲圣教信仰，取代了基督徒。他们不再按宗徒的吩咐去查问：\BibleQuote{若有无可指摘的人}（\BibleRef{铎 1:8}）：却效法那不虔敬的\Person{雅洛贝罕}的做法：谁能提供最多的金钱，谁就被任命为主教；即使那人碰巧是异教徒，只要他向他们提供金钱，这对他们也毫无区别。从\Person{亚历山大}时代起一直担任主教的人、隐修士和\OriginalTerm{苦修者}{ascetics}都被放逐了；而那些只擅长毁谤的人，尽其所能地败坏了宗徒的规范，玷污了教会。的确，他们对我们的诬告为他们赢得了不少，使他们得以在您的时代行此不义，做这样的事；以致于经上的话可以用在他们身上：\BibleQuote{祸哉，那些使我的名在外邦人中受亵渎的人}（\BibleRef{罗 2:24}）。\par

29. \Emph{亚大纳削听闻自己被放逐的消息。}\par

谣言就是如此四处传播的；尽管一切都这样被颠倒了，我仍未放弃渴望来到您的虔敬面前的强烈愿望，而是再次准备启程。我之所以如此急切，是因为我确信这些行为违背了您的意愿，并且如果陛下得知了所发生的事，定会为将来而加以阻止。因为我无法想象，一位正义的君王会希望主教被放逐，贞女被剥去衣服，或者教会受到任何方式的扰乱。当我这样思考并加速行程时，看，又有第三个消息传到我耳中，大意是：已写信给\Place{阿克苏姆}的君主们，要求把阿克苏姆的主教\Person{福鲁门提乌斯}从那里带来，并且甚至到蛮族之地去搜寻我，好把我交给总督们的所谓\Quote{\OriginalTerm{记名卷}{Commentaries}}，并强迫所有平信徒和圣职人员与\OriginalTerm{阿里乌斯异端}{Arian heresy}共融，凡不服从此命令者处死。为了表明这些并非仅仅是空穴来风，而是已被事实证实，既然陛下允许，我便出示这封信。我的敌人们不断宣读它，并以死亡威胁每一个人。\par

30. \Emph{君士坦提乌斯反对亚大纳削的信的副本。}\par

\Person{维克多·君士坦提乌斯·马克西穆斯·奥古斯都}致\Place{亚历山大}人。\par

\Quote{你们的城邑，保持其民族特性，并记念其建立者的德行，素常向我们表示服从，如同今日所行；我们则若不以善意超越\Person{亚历山大}本人，便觉大大亏欠。因为节制之心在凡事上有序，而王权因德行，容我如此说，即如你们这样的德行，当拥抱你们过于万人；你们兴起为首批智慧之师，首先承认天主；更为自己拣选了最成全的导师；且由衷赞同我们的意见，公义地憎恶那骗子与欺诈者，忠信地联合于那些超乎万民景仰的可敬之人。然而，谁人不知，即使住在地极之人，也不知近来所显露的暴烈党争？我们不知何事可与之相比。大多数公民眼睛昏蒙，一个人从最卑污的巢穴出来，在他们中间得了权柄，以虚假诱骗那些渴望认识真理的人——如同在黑暗的掩护下——这人从不给予他们任何有益和启迪的言论，却以无益的诡辩败坏他们的心智。他的谄媚者向他欢呼喝采；他们对他能力惊叹，且可能仍暗自咕哝；而多数单纯的人则效法他们。于是，一切随波逐流，彷佛洪水冲入，万事尽都荒废。众人中有一人掌权——我如何能更真切描述他呢，只能说，他不比最卑微的百姓优越分毫，他给予这城的唯一恩惠，就是没有将市民推下坑中。这位心胸高尚、声名显赫之人，不等审判临身，便自判放逐，恰如其分。现在，蛮族将之除灭，对他们有益，免得他引领他们中一些人陷于不敬；因为他会向首先遇见的人诉苦，如同戏剧中悲苦的角色。然而，我们如今向他永别。对于你们，我鲜能找到可相提并论者：我不得不分别将你们尊于万人之上，因你们的行为所彰显的大德大智，几乎传遍整个世界。继续这清醒的途径吧。我极愿有人以当得的赞美之词对我叙述你们的品行；哦，你们在荣耀的赛跑中超越了前辈，且将作现今在世之人，以及未来世代之人的高贵榜样，唯有你们为自己拣选了最完全之存在，引领你们言行的向导，且在言语和行为上毫不迟疑，勇敢地转移情感，归向另一方，离开那些卑琐尘世的教师，伸展向天上的事物，在最可敬的\Person{乔治}的引导之下，无人比他更受完备的训导。在他的引导下，你们将保持对来生的美善盼望，并在现今世界得享安宁。但愿所有公民一同持守他的言语，如同神圣的锚，如此我们便无需对那些灵魂患病之人动用刀或烧灼！我们极诚恳地劝告这样的人，放弃他们对\Person{亚大纳削}的热忱，甚至不再记念他曾在他们中间大量讲说的愚妄之事。否则，他们终将在不觉间陷于极大危险，我们不知道有谁能巧妙解脱这类党争之人。因为当那个害人虫\Person{亚大纳削}被四处驱逐，因极其卑劣的罪行而被定罪，若将他杀死十次，亦只受当得之刑，我们若容忍那些谄媚者和混淆之徒因我们而欢跃，便不合宜；这类人的品格，提及亦是可耻。关于他们，早已有令交付官长，当将他们处死。但也许现在他们若停息先前的过错，并终能悔改，便不致灭亡。因为那极恶的害人虫\Person{亚大纳削}引领他们，败坏整个邦国，以其不洁和亵渎之手，玷污了至圣之物。}\par

31. \Emph{\Person{君士坦提乌斯}致\Place{埃塞俄比亚}人反对\Person{弗鲁门提乌斯}书}\par

以下是写给\Place{阿克苏姆}的王子们、论及该地主教\Person{弗鲁门提乌斯}的信函。\par

\Person{君士坦提乌斯}常胜至尊奥古斯都，致\Person{埃扎纳}和\Person{萨扎纳}。\par

对我们来说，传播至高天主的认识，是一件至为关切和挂虑的事。我想，全人类在这方面也要求我们给予同等的关怀，好使他们能在希望中度过一生，得以认识天主，并在寻求正义和真理时彼此没有分歧。因此，鉴于你们与\Place{罗马}人一样值得同样的眷顾，并愿对你们的福祉显示同等的关切，我们命令在你们的教会里宣认与他们相同的教义。所以，请尽快派遣主教\Person{弗鲁门修斯}到\Place{埃及}，去见最可敬的\Person{乔治}主教和在那里有着特别权力任命这些职务并裁决相关事宜的其他人。因为你们当然知道并记得（除非你们假装不知道众人皆知的事），这位\Person{弗鲁门修斯}是由犯有万般罪行的\Person{亚大纳削}擢升至现今地位的；\Person{亚大纳削}根本无法公平地洗清对他的任何指控，反而立即被剥夺了主教职位，如今流离失所，漂泊不定，从一个地方漫游到另一个地方，仿佛这样就能逃脱自己的邪恶。现在，如果\Person{弗鲁门修斯}甘心服从我们的命令，并接受对其任命一切情况的审查，他就会向众人清楚表明，他丝毫没有违背教会的法律和公认的信仰。经审讯，他应证明其一贯的良好品行，并向将审判这些事的人陈明自己的生平，若确实显示他有任何权利成为主教，他将从他们那里接受任命。但如果他拖延逃避审讯，那就显然他是受了恶人\Person{亚大纳削}的怂恿，对天主如此放纵不虔敬，甘愿追随那恶行昭彰者的道路。我们担心他会越过边界进入\Place{阿克苏姆}，用可憎而不敬虔的言论败坏你的百姓，不仅扰乱和动荡教会，亵渎至高天主，还会因此使他所造访的各个民族完全倾覆和毁灭。但我确信，\Person{弗鲁门修斯}将返回家乡，完全熟知教会的一切事务，他会从与最可敬的\Person{乔治}和其他那些极有资格传授此种知识的主教们的交谈中，获得大有裨益且普遍有用的许多教导。愿天主永远护佑你们，最尊贵的弟兄们。\par

32. \Emph{他为自己逃跑辩护。}\par

我听见，甚至几乎看见了这些事，通过使者们悲哀的陈述，我承认我又退回旷野，并正确地推断——正如你的虔诚将会觉察到的——如果我被寻找，是为了在被人发现时立刻送到总督那里，我便永远不能来到你的恩宠面前；而那些不愿签署反对我的人遭受如此严重迫害，并且拒绝与亚略派共融的平信徒被下令处死，那么无疑，诽谤者会对我设计出千万种新的毁灭方式；而在我死后，他们会对任何他们想要伤害的人，使用他们选择的任何手段，更加大胆地对我们撒谎，因为那时再也没有人能揭露他们了。我逃跑，不是因为害怕你的虔诚（因为我知道你的宽仁和善良），而是因为从已发生的事中，我看出了仇敌的精神，并考虑到他们会使用一切可能的手段置我于死地，因为他们害怕会被要求为违背你陛下旨意所做的事承担责任。你瞧，你的恩宠只命令将主教们驱逐出城市和省区。但这些\Quote{善良}的人擅自行事，超出你的命令，把年迈且德高望重的主教们放逐到沙漠、无人居住的可怕之地，甚至跨越三个省区的边界。有些人从\Place{利比亚}被送往\Place{大绿洲}；另一些人从\Place{底比斯}被送往\Place{利比亚}的\Place{亚摩尼阿卡}。我逃跑也不是因为怕死；请他们不要指责我犯有懦弱之罪；而是因为我们救主的命令：在受迫害时要逃跑，在被迫寻时要隐藏，不把自己暴露在必然的危险中，也不在迫害者面前出现而更激起他们对我们的愤怒。因为把自己交给敌人去杀害，就等于杀害自己；但按照救主的命令逃跑，意味着认清时机，并表现出对迫害者真正的关怀，免得他们若流人血，就会触犯\BibleQuote{不可杀人}这条律法\BibleRef{出 20:13}。然而这些人却用诽谤我的言辞，急切希望我被害。他们最近再次做出的事证明这就是他们的愿望和杀人的意图。我相信，至爱于天主的\Person{奥古斯都}，你听了一定会惊愕；这真是一种令人震惊的暴行。我恳请你简短地听我道来。\par

33. \Emph{亚略派对献身贞女的行为。}\par

圣子，我们的主和救主\Person{耶稣基督}，为了我们而降生成人，毁灭了死亡，并将我们的种族从败坏的奴役中解救出来，此外，在祂所有其他恩惠之外，祂还赐给我们这一恩惠：我们应在世上，于\OriginalTerm{童贞}{virginity}的状态中，拥有天使圣洁的肖像。因此，\OriginalTerm{天主教会}{Catholic Church}惯常称那些达到此德行者，为\OriginalTerm{基督的净配}{brides of Christ}。而看见她们的异教徒，则赞叹她们为\OriginalTerm{圣言的圣殿}{temples of the Word}。因为确实，这神圣而属天的职业，仅在我们基督徒中建立，且是一个极有力的证明，表明在我们中间能找到纯正真实的宗教。你那位最虔诚的父亲、受福忆念的奥古斯都\Person{君士坦丁}，特别尊崇童贞女超过一切人，而陛下在若干信函中，赐予了她们\Quote{可敬而圣洁的女性}的称号。但如今，这些诽谤我的可悲\OriginalTerm{亚略派}{Arians}——他们曾阴谋反对大多数主教——取得了长官的同意与合作，先剥去她们的衣服，然后把她们悬吊在所谓的\OriginalTerm{赫尔墨塔里刑架}{Hermetaries}上，再三猛烈鞭笞她们的肋部，连真正的罪犯都未曾遭受此等酷刑。\Person{彼拉多}为了讨好昔日的犹太人，用长矛刺穿了我们救主的一侧；这些人比彼拉多更为疯狂，因为他们鞭打了祂身体的两侧，而非一侧；因为童贞女的肢体，在特殊意义上是属于救主的。所有人听到这类行径的简单叙述都不寒而栗。惟独这些人不但不怕剥去并鞭笞那些童贞女已全然奉献于我们救主基督的无玷肢体；更恶劣的是，当众人因如此极端的残酷而谴责他们时，他们丝毫不以为耻，反而伪称这是出于陛下的命令。如此，他们胆大妄为，满心恶念与意图。这类暴行在过去的教难中闻所未闻；即便从前曾发生过，在陛下——一位基督徒——的治下，童贞竟遭受如此凌辱与羞耻，或这些人竟将自身的残酷归咎于陛下，这绝不相宜。这种邪恶只配异端者所为，即亵渎圣子，并向祂的圣洁童贞女施暴。\par

34. \Emph{他规劝\Person{君士坦提乌斯}}\par

当\OriginalTerm{亚略派}{Arians}再次犯下如此暴行时，我遵从圣经的指示，\BibleQuote{你藏身片时，等到主的忿怒过去}，这确实没有错。这是我退隐的另一原因，最蒙天主喜爱的奥古斯都；我不拒绝退入旷野，也不拒绝在必要时从城墙上用篮子缒下。\BibleRef{格后 11:33} 我忍受一切，甚至住在野兽之中，为使你的恩宠可重归于我，等候时机向你呈上此辩护，深信他们必将被判罪，而你的良善必将向我彰显。哦，奥古斯都，受福且最蒙天主喜爱的，你要我怎样做？要我在诬告者对我充满怒火、意图杀我之时来见你？还是如经上所记，\BibleQuote{暂时藏身}，使他们在此期间，因是异端者而被定罪，你的良善得以向我彰显？或者，陛下，你要我出现在你的长官面前，以致你虽仅以威吓的言辞写信，他们却不领会你的用意，反而因亚略派而迁怒于我，凭你的书信而杀我，并将谋杀归咎于你？无论是我屈服而任人流血，或你身为基督徒君王，竟有杀害基督徒——且是主教们——的罪名归咎于你，都绝不相宜。\par

35. 因此我藏身并等候此时机，实为更佳。是的，我确信，从你对圣经的认识，你必会同意并赞同我在此事上的行为。因为你将看到，如今那些激怒你反对我们的人已被静默，你正直的仁慈彰显出来，并向众人证明，你从不迫害基督徒，而是他们使教会荒凉，以便到处播下自己不敬的种子；为此，若非我逃亡，我早就遭受他们的毒手。因为显然，那些不惮于在伟大的奥古斯都前如此诽谤我，并如此猛烈攻击主教和童贞女的人，也企图谋害我的性命。但感谢主，祂将王国赐给了你。众人都坚定了对你良善的看法，也坚定了对他们邪恶的看法；我最初逃离这些，为能如今向你发出此呼吁，并让你能找到可施恩的对象。我恳求你，因此，为经上记着说：\BibleQuote{温和的回答平息忿怒}，以及\BibleQuote{正直的思想为君王所悦纳}；请接受我这辩护，并使所有主教及其余神职人员回归他们的家乡和教会；这样，控告我之人的邪恶便显露出来，而你，无论现今还是在审判之日，都能坦然向我们的主和救主、万王之王耶稣基督说：\Quote{\BibleQuote{你所交给我的人，我连一个也没有失落，}} \BibleRef{若 18:9}但这些人却是图谋毁灭所有的人，我为那些丧亡者悲伤，为被鞭笞的童贞女，并为所有针对基督徒的其他恶行悲伤；我将被放逐者领回，并使他们归返各自的教会。\par

\SourceRecord{Translated by M. Atkinson and Archibald Robertson.；Nicene and Post-Nicene Fathers, Second Series；Vol. 4.；Edited by Philip Schaff and Henry Wace.；Buffalo, NY: Christian Literature Publishing Co.,；1892.；<http://www.newadvent.org/fathers/2813.htm>.}
